% F03：本书原创矢量示意；依据所在节的定义、证明或明确模型。
\begin{figure}[htbp]
\centering
\begin{tikzpicture}[x=1cm,y=1cm,>=stealth,every node/.style={font=\small},line width=.65pt]

\draw[->,gray] (0,1.5)--(5.3,1.5) node[right]{\(x\)};
\fill[AnalysisBlue!15] (0,1.5)--(0,.4)--(2.2,1.5)--cycle;
\fill[orange!20] (2.2,1.5)--(4.4,2.6)--(4.4,1.5)--cycle;
\draw[thick,AnalysisBlue] (0,.4)--(4.4,2.6);
\draw[->,very thick,orange!80!black] (0,1.5)--(0,3);
\node at (2.2,3.5) {\(\nu=fm_1+\delta_0\)};
\node at (2.2,-.05) {带符号密度与原子};
\draw[->,gray] (7,0)--(12.3,0) node[right]{\(x\)};
\fill[AnalysisBlue!15] (7,0)--(7,1.1)--(9.2,0)--(11.4,1.1)--(11.4,0)--cycle;
\draw[thick,AnalysisBlue] (7,1.1)--(9.2,0)--(11.4,1.1);
\draw[->,very thick,orange!80!black] (7,0)--(7,2.2);
\node at (9.2,3.5) {\(|\nu|=|f|m_1+\delta_0\)};
\node at (9.2,-.5) {全变差累加两侧质量};

\end{tikzpicture}
\caption[符号测度的正负质量与全变差]{符号测度的正负质量与全变差。这里 \(f=(2x-1)\cdot\mathbf1_{[0,1]}\)。折线表示连续部分的密度，箭头表示单位原子。箭头高度仅用于辨认原子，不与密度纵坐标作面积比较。}
\label{b1:fig:visual-f03}
\end{figure}
